% TOP_PROBLEM: {"id": 3700017, "title": "Levin's derivative-zero problem", "category_id": 8, "category": {"id": 8, "name": "analysis", "display_name": "Analysis", "description": "Limits, continuity, calculus, and function theory.", "slug": "analysis", "order_index": 8, "created_at": "2026-07-31T15:26:25.671Z"}, "status": "open", "problem_number": "AMR-036-0017", "set_id": 13, "statusQualification": "Uses the author's 10 September 2026 revision c(exp(ibz)-a), with arbitrary complex a, rather than the superseded unit-modulus subtraction in the catalogue. Identically zero derivatives are excluded from the zero test."}
% FIRST_POSTED: 2026-10-01
% ENTRY_KIND: statement_only
% UnsolvedMath ID 3700017; source catalogue code AMR-036-0017
% !TeX program = lualatex
% Definition and sources only; this file is not an LLM solution attempt.
\documentclass[11pt,a4paper]{article}
\usepackage[margin=25mm]{geometry}
\usepackage{fontspec}
\setmainfont{lmroman10-regular.otf}[BoldFont=lmroman10-bold.otf,
 ItalicFont=lmroman10-italic.otf,BoldItalicFont=lmroman10-bolditalic.otf]
\usepackage{amsmath,amssymb,mathtools}
\usepackage{unicode-math}
\setmathfont{latinmodern-math.otf}
\AtBeginDocument{\let\setminus\smallsetminus}
\usepackage{xurl,hyperref}
\hypersetup{colorlinks=true,urlcolor=blue}
\setcounter{secnumdepth}{0}
\setlength{\parindent}{0pt}
\setlength{\parskip}{5pt}
\setlength{\emergencystretch}{3em}
\begin{document}
\section*{Levin's derivative-zero problem}\label{3700017}
{\small UnsolvedMath ID 3700017 · AMR-036-0017 · Analysis · AMR Open Problem Lists}
\subsection{Definitions and mathematical statement}
Let $\mathbb H_-=\{z\in\mathbb C:\operatorname{Im}z\leq0\}$.
Suppose $f$ is a nonzero entire function such that every
zero of every nonidentically-zero derivative $f^{(n)}$,
$n=0,1,2,\ldots$, belongs to $\mathbb H_-$.
Identically zero derivatives are disregarded, so that
polynomials are included in the natural way.

Let $\mathcal L_-$ be the compact-open closure of the
polynomials whose zeros all lie in $\mathbb H_-$: membership
means being a locally uniform limit on $\mathbb C$ of
such polynomials. Classify all entire functions satisfying
the derivative-zero condition. In particular, must every
one belong either to $\mathcal L_-$ or to one of the
exceptional families
\[
             f(z)=c e^{az},\qquad
             f(z)=c(e^{ibz}-a),
 \quad c,a\in\mathbb C,\quad b\in\mathbb R?
\]
Only members of these families satisfying the original
zero-location hypothesis are relevant.

This uses the broadened exceptional family in the author's
September 2026 revision: the subtracted constant is allowed
to be complex, rather than restricted to the unit circle.
More generally, the classification target is to describe
precisely any exceptional family outside $\mathcal L_-$
if the displayed list is incomplete.
The half-plane is fixed for all derivative orders. Separate
half-planes chosen for different derivatives, or conditions
on only finitely many derivatives, would give weaker
problems.
\subsection{Short English statement}
Classify entire functions whose zeros and all derivative zeros stay in one closed half-plane.
\subsection{Sources}
\begin{itemize}
\item[S1] ULAM.AI and the UnsolvedMath Contributors, supplied problems.json, record 3700017 (AMR-036-0017), AMR Open Problem Lists. Catalogue identity and source transcription; CC BY 4.0.
\url{https://huggingface.co/datasets/ulamai/UnsolvedMath}.
\item[S2] Eremenko - My favorite unsolved problems, Problem of B. Ya. Levin. Original source indexed by AMR-036-0017.
\url{https://www.math.purdue.edu/~eremenko/uns1.html}.
\item[S3] A. Eremenko, A problem of B. Ya. Levin. Author-maintained primary problem note.
\url{https://www.math.purdue.edu/~eremenko/dvi/levin.pdf}.
\end{itemize}
\end{document}
